\subsection*{Choosing a successor absent from the visible history}
Two extracts from the same article can end in the same character and continue
with different ones. In 512 held-out Chinese articles, we supplied two such
histories and both possible successor images (Fig.~\ref{fig:visual-successors}).
Neither successor appeared in either history. The final native writer assigned
the higher score to the same candidate in both contexts for 509 articles. It
selected both successors independently in only two articles. A joint assignment
can detect smaller relative shifts in preference. The final writer made 280/512
correct joint assignments, compared with 281/512 at initialization.

A separate visual reader learned a compatibility score between preceding ink
and a proposed successor. Training on 1,047,284 glyph transitions across 2,127
modern passages took 19.46 seconds on two CPU threads. The reader assigned
282/512 pairs correctly and chose both independent successors in 45 articles.
Reversing the first 31 cells while retaining the final cell reduced its joint
score by 8.2 percentage points (95\% article bootstrap interval 2.3 to 13.9).
Its 4.3-point advantage over rotated supervision had an interval of $-2.0$ to
10.7 points. Its 1.8-point advantage over initialization had an interval of
$-4.5$ to 8.0. The prespecified training criterion therefore remained unmet.

An otherwise matched character-indexed reader scored 288/512 after 18.73
seconds. The visual reader used 121,312 parameters and the character-indexed
reader 584,448. The visual minus character-indexed accuracy difference was
$-1.2$ points (interval $-6.8$ to 4.5). Counts of three consecutive
physical glyphs achieved 62.7\% joint assignment accuracy with half
credit for 288 ties. Time to a common level of independent language
generation quality remains unmeasured.
